\documentclass[12pt,leqno]{article}
\usepackage{amsmath,amssymb,amsthm,colortbl,xcolor,nicefrac,setspace}
\usepackage[T1]{fontenc}
\usepackage[expert,seriftt,lucidasmallscale]{lucidabr}
\usepackage{mathrsfs,hyperref,color,graphicx}
\hypersetup{hidelinks}
\setlength{\textwidth}{6.5in}
\setlength{\textheight}{9in}
\setlength{\oddsidemargin}{0mm}
\setlength{\evensidemargin}{0mm}
\setlength{\topmargin}{0mm}
\setlength{\headheight}{0mm}
\setlength{\headsep}{0mm}
\setlength{\topskip}{0mm}
\setlength{\hoffset}{0in}
\setlength{\voffset}{0in}
\newcommand{\AND}{\ensuremath{\mathbin{\&}}}
\newcommand{\OR}{\ensuremath{\vee}}
\newcommand{\IFF}{\ensuremath{\leftrightarrow}}
\newcommand{\IC}{\ensuremath{\rightarrow}}
\newcommand{\IF}{\ensuremath{\supset}}
\newcommand{\NOT}{\ensuremath{\text{$\thicksim$}}}
\newcommand{\U}{\mathcal U}
\newcommand{\two}{\{0,1\}}
\newcommand{\Atom}{\operatorname{Atom}}
\onehalfspacing
\begin{document}
\begin{center}
{\large A proposed counterexample to Proposition 2.11}\par
\small \emph{Classicism}, 1 July 2022 draft, p.~30 and footnote~42
\end{center}

The following construction appears to refute Proposition 2.11, not just
the normalization in footnote~42. It satisfies Classicism and, at every
world, BF, Boolean Completeness, and Atomicity. Nevertheless, Rigid
Comprehension fails for properties of propositions.

The construction uses a nonprincipal ultrafilter, but only ultralimits
in fixed finite sets. We give the carriers and action maps, verify
closure and the hypotheses, and exhibit the failed comprehension instance.

\medskip
\noindent\textbf{1. The typed model.}
Work in ZFC with a nonprincipal ultrafilter $\U$ on the positive integers.
Take one full extensional simple-type model $M$, with $M_e=\{*\}$,
$M_t=\two$, and full function spaces. Each fixed type carrier is finite.
For a sequence $(a_n)$ in a finite nonempty set $A$, define
$\lim_\U a_n=a$ when $\{n:a_n=a\}\in\U$. Exactly one cell of a finite
partition belongs to an ultrafilter, so this value exists and is unique.
If $f_n:A\to B$ and $a_n\in A$, with $A,B$ finite, then
\begin{equation}\label{eq:limit}
 \lim_\U f_n(a_n)=(\lim_\U f_n)(\lim_\U a_n).
\end{equation}
Intersect the two $\U$-large sets on which $f_n$ and $a_n$ take their
limiting values to obtain this identity. The same argument applies to
each operation on fixed finite sets. Changing finitely many terms does
not change a limit, since $\U$ contains every cofinite set.

Use worlds and arrows
\[
o\longrightarrow s\longrightarrow u,\qquad
o\longrightarrow n\quad(n\geq1),
\]
including identities and the composite $o\to u$. Put $M$ at every $n$
and at $u$. At $s$, take $D_s(e)=\{*\}$ and $D_s(t)=\two^2$, with
$j_e=\mathrm{id}$ and $j_t(b_s,b_u)=b_u$. Recursively, $D_s(\alpha\to\beta)$
contains all pairs $(g_s,g_u)$ with
$g_s:D_s(\alpha)\to D_s(\beta)$, $g_u\in M_{\alpha\to\beta}$, and
$j_\beta g_s=g_u j_\alpha$. Define
$\operatorname{app}_s((g_s,g_u),x)=g_s(x)$; the action
$j_{\alpha\to\beta}$ returns $g_u$.

At $o$, take $D_o(e)=\{*\}$ and
\[
 D_o(t)=\two\times\two\times\two^{\mathbb N_{>0}},\qquad
 \pi_s(b_o,b_s,(b_n))=(b_s,\lim_\U b_n),\quad \pi_n(b)=b_n.
\]
A proposition is true at its current world exactly when its current bit
is 1. A root object of type $\alpha\to\beta$ is a triple
$(f_o,F_s,(f_n))$, where
\[
 f_o:D_o(\alpha)\to D_o(\beta),\quad
 F_s\in D_s(\alpha\to\beta),\quad f_n\in M_{\alpha\to\beta}.
\]
For every $x\in D_o(\alpha)$ it satisfies
\[
 \pi_s(f_o(x))=\operatorname{app}_s(F_s,\pi_s x),\quad
 \pi_n(f_o(x))=f_n(\pi_n x),\quad j(F_s)=\lim_\U f_n.
\]
Application uses the current function component. The maps to $s$ and $n$
return the corresponding future components; the map to $u$ is $j\pi_s$.


\medskip
\noindent\textbf{2. Surjectivity and application.}
We prove by induction on types that every carrier is nonempty,
each $D_s(\alpha)$ is finite, $j_\alpha$ is onto, and
\[
 x\longmapsto(\pi_sx,(\pi_nx)_n)
\]
maps $D_o(\alpha)$ onto
\[
 E_\alpha=
 \{(a_s,(a_n)):a_s\in D_s(\alpha),\ a_n\in M_\alpha,\
                         j_\alpha a_s=\lim_\U a_n\}.
\]
For $e$ these claims follow from the singleton definitions. For $t$,
a compatible pair consists of $a_s=(b_s,b_u)$ and a sequence $(b_n)$
with $b_u=\lim_\U b_n$; the root proposition $(0,b_s,(b_n))$ lifts it.
The map $j_t$ is onto.

Suppose the assertions hold for $\alpha,\beta$.
The intermediate carrier at $\alpha\to\beta$ is a subset of a product
of finite function spaces. Given $g_u\in M_{\alpha\to\beta}$, choose,
for each $x\in D_s(\alpha)$, a preimage under $j_\beta$ of
$g_u(j_\alpha x)$. These choices define $g_s$, so $(g_s,g_u)$
belongs to $D_s(\alpha\to\beta)$ and maps to $g_u$.

For root joint surjectivity, let $(F_s,(f_n))\in E_{\alpha\to\beta}$.
For each $x\in D_o(\alpha)$ the desired outputs are
\[
 y_s=\operatorname{app}_s(F_s,\pi_sx),\qquad y_n=f_n(\pi_nx).
\]
They are compatible, since
\[
 j_\beta y_s
 =j_{\alpha\to\beta}(F_s)(j_\alpha\pi_sx)
 =(\lim_\U f_n)(\lim_\U\pi_nx)
 =\lim_\U y_n.
\]
The induction hypothesis at $\beta$ supplies a root preimage of this
tuple. Choose one for each $x$; these choices give the required $f_o$.
This uses ordinary metatheoretic choice in ZFC. Nonemptiness follows
at the same step by taking constant functions between inhabited carriers.

Every individual root action is onto. To lift $a_s$, take the constant
sequence $a_n=j_\alpha a_s$ and use joint surjectivity. To prescribe
a value at one terminal world $n$, start with such a constant sequence
and change its $n$th entry. The limit is unchanged. The action to $u$
is onto because it is the composite $j_\alpha\pi_s$.
Application commutes with every action by the defining equations.

The current application function determines each function object.
At $s$, surjectivity of $j_\alpha$ determines $g_u$ from $g_s$.
At $o$, surjectivity of $\pi_s$ and each $\pi_n$ determines $F_s$
and $f_n$ from $f_o$, using the intermediate result for $F_s$.
Equality here is equality of whole objects, not mere agreement of
their current truth extensions.

\medskip
\noindent\textbf{3. Logical operations and abstraction.}
For an arrow $h:v\to w$, write $h_\alpha$ for its action at type
$\alpha$. The truth bit at $w$ of equality at $v$ is
\[
 [\,h_\alpha a=h_\alpha b\,],
\]
where brackets denote $1$ if the statement is true and $0$ otherwise.
Negation, conjunction and disjunction act on the truth bits at each
world. For $F\in D_v(\alpha\to t)$, the truth bit at $w$ of
$\forall_\alpha F$ is
\[
 [\,\operatorname{app}_w(h_{\alpha\to t}F,a)
       \text{ is true at }w\text{ for every }a\in D_w(\alpha)\,].
\]
Existential quantification uses \emph{some} in place of \emph{every}.

These prescriptions define admissible objects. At $s$, their future
components are the corresponding operations in $M$. At the root,
the required $u$-bit is the ultralimit of the terminal bits.
For equality this follows from the finite-set limit identity applied
to the equality test. For quantification, the terminal predicate
functions range over the finite set $M_{\alpha\to t}$; on a
$\U$-large set the entire function equals its limit. Universal and
existential quantification over the fixed finite $M_\alpha$ therefore
give the same limiting bit. The Boolean cases follow from
\eqref{eq:limit}. Their curried function components at future worlds
are specified by the same prescriptions, so application naturality and
the higher-type compatibility equations hold as well.

Interpret a variable by its assigned object, application by the current
application function, and a constant at a future world by transporting
its root value. At $v$, interpret $\lambda x.A$ by its current function
\[
 a\longmapsto\lbrack\!\lbrack A\rbrack\!\rbrack_v^{g[x\mapsto a]}
\]
together with the corresponding future functions. Simultaneous
induction on terms proves membership in the prescribed carriers and
\begin{equation}\label{eq:naturality}
 h_\alpha\lbrack\!\lbrack A\rbrack\!\rbrack_v^g
       =\lbrack\!\lbrack A\rbrack\!\rbrack_w^{hg}
       \qquad(A:\alpha,\ h:v\to w).
\end{equation}
For abstraction, apply the induction hypothesis for the body
to $g[x\mapsto a]$; its transported assignment is
$(hg)[x\mapsto h_\sigma a]$, where $x:\sigma$. This is the
compatibility equation for the current and future functions.
Application uses its defining naturality equation, and logical
constants use the preceding paragraph.

At the root one must also check the ultralimit condition between
the terminal functions and the $u$-function. Each term involves only
finitely many free variables and nonlogical constants. Intersect the
$\U$-large sets on which their terminal values equal their $u$-values.
On this intersection the interpretations in the identical full model
$M$ agree. This proves the limit condition for the whole denotation,
including function-valued terms. Thus the abstraction functions above
belong to the required carriers.

For comparison with Definitions 3.18--3.20, encode $p\in D_v(t)$
as the set of outgoing arrows $h:v\to w$ for which $h_t p$ is
true at $w$. Recursively encode a function $F$ by the graph
\[
 (h,a)\longmapsto\operatorname{app}_w(h_{\alpha\to\beta}F,a)
       \qquad(a\in D_w(\alpha)).
\]
Arguments and outputs are encoded at their respective types.
The proposition encoding is injective: it records every free bit.
The function encoding is injective by induction, since it includes
current application and all future application. Transport becomes
the paper's inverse-image action on propositions and precomposition
action on graphs. The truth clauses above are its six primitive
clauses, and \eqref{eq:naturality} supplies its abstraction clause.
We therefore obtain an action model in the paper's sense.
Theorem~3.23 gives Classicism.

In this model $\Box p$, defined as $p=\top$, is true at $v$ just when
$p$ is true at every world accessible from $v$: equality with $\top$
compares the whole outgoing truth profile.

\medskip
\noindent\textbf{4. BF, Boolean Completeness, and Atomicity.}
For BF, suppose $\forall x\Box P(x)$ holds at $v$.
Given $h:v\to w$ and $a\in D_w(\alpha)$, choose $b\in D_v(\alpha)$
with $h_\alpha b=a$. The assumption at $b$ gives $P(a)$ at $w$,
with the other parameters transported along $h$. Since $a$ and $h$
were arbitrary, $\Box\forall xP(x)$ holds at $v$. This is the
argument in Proposition~3.24(ii); it works at every world and type.

For a relational type $\tau=\sigma_1\to\cdots\to\sigma_k\to t$, put
$T_v=\prod_iD_v(\sigma_i)$. Simultaneous induction on the curried result
type gives
\[
 M_\tau\cong 2^{T_u},\qquad D_s(\tau)\cong 2^{T_s}\times2^{T_u},
\]
\[
 D_o(\tau)\cong 2^{T_o}\times
 \{(S,(N_n)):S\in D_s(\tau),\ N_n\in M_\tau,\ jS=\lim_\U N_n\}.
\]
The case $k=0$ is the specified proposition carrier. At the arrow step,
choose a current truth relation and compatible future operators. For each
current argument, the result-type induction supplies the unique output
with that truth section and those future components; ultralimit
commutation with application gives compatibility. Conversely, application
recovers these data uniquely. This also shows that the current relation
at $s$ is independent of its terminal operator.

All $T_n$ equal the fixed finite $T_u$. Thus $(N_n)$ determines the
terminal component of $S$, leaving its current relation free, and
\[
 D_o(\tau)\cong 2^{T_o}\times2^{T_s}\times\prod_{n\geq1}2^{T_n}.
\]
The Boolean operations act coordinatewise. The source order
$A\leq B$, meaning $B=A\OR B$ as whole objects, is therefore the product
order. The root algebra is complete and atomic; at the other worlds the
corresponding algebras are finite. Each factor is the powerset algebra of its tuple
set. Arbitrary meets and joins are intersections and unions in each
free coordinate. Each nonzero element contains a tuple in one of
those coordinates; the element supported on that one tuple in that
one coordinate is an atom below it. The empty tuple is used when
$k=0$, so this includes propositions.

To check the source formulas, let $Q\in D_v(\tau\to t)$ and form
$S_Q=\{a\in D_v(\tau):Qa\text{ is true at }v\}$.
Its internal meet $m$ belongs to $D_v(\tau)$. For each $z$ in that
carrier,
\[
  (\text{for every }a\in S_Q,\ z\leq a)
        \quad\Longleftrightarrow\quad z\leq m.
\]
This is the $\mathrm{GLB}_\tau(m,Q)$ clause, including
when $S_Q$ is empty. The source formula
\[
 \mathrm{Atom}_\tau(y)\ \equiv\
 \forall z\bigl((z\leq y\AND z\neq y)\IFF z\leq\NOT_\tau z\bigr)
\]
says that the strict lower elements of $y$ are precisely bottom.
It excludes $y=\bot$ and holds exactly at the ordinary nonzero
atoms. Hence Boolean Completeness and Atomicity hold at every
world and every relational type, as do their boxed forms and
$\Box\mathrm{BF}$. A GLB is chosen separately at the world where
its existence is asserted. Neither Boolean Completeness nor its
necessitation requires the transport of that GLB to be the GLB
of the transported family.

\medskip
\noindent\textbf{5. The normalization loses its witness.}
Put $w=(0,1,(0)_n)$ and $q=(0,0,(1)_n)$.
The proposition $w$ is a root Boolean atom selecting $s$ (not a true
root-world proposition). Since $\pi_s(q)=(0,1)$ is an atom of $\two^2$,
$w\leq\Atom(q)$. But the actual extension of
$F(p):=(w\leq(p=q))$ is
\[
 \{p:p_s=0\ \text{and}\ \{n:p_n=1\}\in\U\}.
\]
Its meet is $\bot$: the root bit can be zero, the $s$-bit is always zero,
and each fixed $n$-bit can be zero in a member, using a cofinite set
omitting $n$. Every member maps at $s$ to $(0,1)$, whereas its meet maps
to $(0,0)$. Thus $w\not\leq(q=\bigwedge F)$, and the replacement is not
possibly atomic. Completeness of both algebras does not make the
surjective Boolean homomorphism $\pi_s$ preserve infinite meets.
Indeed, the ambient intersection of the equality fiber, as sets of
outgoing arrows, is $\{o\to u\}$. This is not a root proposition:
zero terminal bits force a zero $u$-bit. Its greatest available lower
bound is empty. The internal meet is not the ambient intersection.

\medskip
\noindent\textbf{6. Rigid Comprehension itself fails.}
For a unary property the source definitions are
\begin{align*}
 \mathrm{Persistent}(Z)&\ \equiv\ Z\leq(\lambda x.\Box Zx),\\
 \mathrm{Inextensible}(Z)&\ \equiv\
 \Box\forall Y\bigl(\forall x(Zx\IC\Box Yx)\IC Z\leq Y\bigr).
\end{align*}
A rigid property satisfies both. Persistence says that at every
accessible world, each instance remains an instance along every
further arrow. The order $Z\leq Y$ at $v$ means inclusion of their
extensions at every world accessible from $v$, not just at $v$.
Write $Z_v$ for the current extension of the transported property.

Write $a_0=(0,0)\in D_s(t)$ and set
$A_0=\pi_s^{-1}(\{a_0\})$.
Its $s$-image is $\{a_0\}$ and its $u$-image is $\{0\}$, but each
$n$-image is all of $\two$: a sequence supported at a single $n$ still
has ultralimit zero. Choose $X:t\to t$ with root extension $A_0$ and
all future extensions empty. The displayed Boolean decomposition
admits this property.

Suppose that a rigid property $Z$ is coextensive with $X$ at $o$.
Persistence forces $Z_n=\two$ for every $n$ and $a_0\in Z_s$.
Carrier compatibility then forces $Z_u=\two$: its extension is the
ultralimit of the $Z_n$, not the image along an arrow from a leaf to $u$.
Now choose $B:t\to t$ with extensions
\[
 B_o=A_0,\qquad B_s=\{a_0\},\qquad B_n=B_u=\two.
\]
Every actual $Z$-instance necessarily satisfies $B$. Root
inextensibility gives $Z\leq B$, hence $Z_s\subseteq\{a_0\}$;
persistence gives the reverse inclusion.
At $s$, however, the haecceity $Y$ of $a_0$ has extensions
$Y_s=\{a_0\}$ and $Y_u=\{0\}$.
Every $s$-instance of $Z$ necessarily satisfies $Y$, but $Z\leq Y$
fails at $s$: the element $1\in Z_u$ is not in $Y_u$.
This contradicts the $s$-instance of the \emph{outer-boxed}
inextensibility condition. No rigid property is root-coextensive with
$X$. The individual sort can remain a singleton throughout.
The outer box is essential: $B$ satisfies persistence and the unboxed
inextensibility clause at $o$, but fails source-defined inextensibility
because that clause must also hold at $s$.

\medskip
\noindent\textbf{7. The full setup of footnote~42 is realized.}
For $a\in A_0$, its haecceity has extension $\{a\}$ at $o$,
$\{a_0\}$ at $s$, and $\{\pi_n a\}$ at each $n$.
Let $X^*$ be the actual LUB of these haecceities.
The root Boolean join is coordinatewise in the $o,s,n$ extensions;
its $u$-extension is then determined by the ultrafilter. This gives
\[
 X^*_o=A_0,\qquad X^*_s=\{a_0\},\qquad X^*_n=X^*_u=\two.
\]
Take $z=q$ as a proposition-input object, and take a root property $Y$
with $Y_o=\varnothing$, $Y_s=\{a_0\}$, and $Y_n=Y_u=\{0\}$.
At $s$, every $X^*$-instance necessarily satisfies $Y$, while
$X^*(z)\AND\NOT Y(z)$ has $s/u$-profile $(0,1)$. Thus
\[
 w\leq\bigl[\forall x(X^*(x)\IC\Box Y(x))\AND\Atom(q)
        \AND q\leq(X^*(z)\AND\NOT Y(z))\bigr].
\]
These are the preceding witness conditions, with $q$ playing the role
of the footnote's $w'$. Its equality-fiber meet is still bottom.

\medskip
\noindent\textbf{8. A citation.}
Footnote~42(ii) should cite Proposition~2.7 rather than~2.6,
which presupposes $C5$. This does not supply the meet-preservation
step at issue.

\medskip
So the issue seems to be the assumed preservation of this meet, rather
than a missing necessity operator on Boolean Completeness. The distinction is between existence of internal meets and their
preservation under the action maps. The construction has the former,
even at every world, but not the latter.
\end{document}
